\appendix

% \section{Inserted Channel Evaluation Data}
% \label{app:a}

% \begin{figure*}[h]

%     \centering
%     \includegraphics[width=\linewidth]{figures/PLA_insert_eval.pdf}
%     \caption{Light transmission through PLA printed channels which were inserted into the holder afterwards. Each subfigure shows the photodiode readings for red and blue light across varying parameters (a) length, (b) curvature, and (c) thickness. No significant differences were observed across parameters.}
%      \Description{} 
%     \label{fig:pla_insert}
%     \vspace{-2mm}
% \end{figure*}
% \begin{figure*}[h]

%     \centering
%     \includegraphics[width=\linewidth]{figures/PETG_insert_eval.pdf}
%     \caption{Light transmission through PETG printed channels which were inserted into the holder afterwards. Each subfigure shows the photodiode readings for red and blue light across varying parameters (a) length, (b) curvature, and (c) thickness. Inserted PETG channles also show no significant difference.}
%      \Description{} 
%     \label{fig:petg_insert}
%     \vspace{-2mm}
% \end{figure*}
% \begin{figure*}[h]

%     \centering
%     \includegraphics[width=\linewidth]{figures/TPU_insert_eval.pdf}
%     \caption{Light transmission through TPU printed channels which were inserted into the holder afterwards. Each subfigure shows the photodiode readings for red and blue light across varying parameters (a) length, (b) curvature, and (c) thickness. Inserted TPU channels show no significant difference.}
%      \Description{} 
%     \label{fig:tpu_insert}
%     \vspace{-2mm}
% \end{figure*}

% For PLA, PETG, and TPU, the inserted channels showed minimal loss, with photodiode readings remaining relatively constant near the maximum for channel length, curvature, and thickness as seen in Figures \ref{fig:pla_insert}, \ref{fig:petg_insert}, and \ref{fig:tpu_insert}.


\section{Circuit Diagram}
\label{app:ckt}
In our experiments, the photodiode was connected in a reverse-bias configuration with a load resistor to form a transimpedance circuit. The cathode of the photodiode was tied to the 5 V supply, and the anode was connected to the Arduino analog input (A0). A 1 M$\Omega$ resistor was placed from the analog input to ground, so that the photocurrent generated under illumination flowed through the resistor, developing a voltage that the Arduino’s ADC could measure.
%In our experiments, the photodiode was connected as a voltage-divider circuit. One terminal of the photodiode was tied to the supply, and the other terminal was connected to an analog input on the Arduino (A0). A load resistor (1 M$\Omega$) was placed from the analog input node to ground. This forms a divider: the photocurrent generated by the diode under illumination flows through the load resistor, developing a voltage that is sampled by the microcontroller’s ADC.
    
    \begin{figure}[h]
    \centering
\includegraphics[width=0.9\linewidth]{figures/circuit2.png}
    \caption{Diagram of circuit used for evaluations.} 
    \Description{the photodiode was connected in a reverse-bias configuration with a load resistor to form a transimpedance circuit. The cathode of the photodiode was tied to the 5 V supply, and the anode was connected to the Arduino analog input. A 1 M$\Omega$ resistor was placed from the analog input to ground, so that the photocurrent generated under illumination flowed through the resistor, developing a voltage that the Arduino’s ADC could measure.} 
    \label{fig:circuit}
    \vspace{-3mm}
\end{figure}


\section{Printing Parameters}
\label{app:params}
Table \ref{tab:printing_params} outlines the parameters utilized for printing length, thickness and curvature evaluation samples.

\begin{table}[h]
\centering
\caption{3D printing parameters used for PLA \& TPU samples for length, thickness, and curvature evaluations.}
\Description{Three columns: Parameter, PLA, and TPU. Each parameter and the corresponding values are written in the columns.}
\label{tab:printing_params}
\begin{tabular}{lcc}
\hline
\textbf{Parameter}            & \textbf{PLA}                 & \textbf{TPU} \\ \hline
Nozzle diameter               & 0.4 mm                       & 0.4 mm \\
Layer height                  & 0.4 mm                       & 0.4 mm \\
Line width                  & 0.41 mm                       & 0.41mm \\
First layer height            & 0.4 mm                      & 0.4 mm \\
Nozzle temperature            & 220\,$^{\circ}$C             & 228\,$^{\circ}$C \\
Bed temperature               & 55\,$^{\circ}$C              & 60\,$^{\circ}$C \\
Infill density (channel)                & 100\% (concentric) & 100\% (concentric ) \\
Infill density (part)                & 15\% (grid) & 50\% (lines ) \\
Cooling fan                   & 100\% after 1st layer        & 100\% after 1st layer \\
\hline
\end{tabular}
\end{table}



\begin{figure*}[ht]
    \centering
    \begin{tikzpicture}[node distance=2mm]
    \node (s11) {\includegraphics[width=0.30\linewidth]{figures/steps_screenshot/1.jpg}};
    \node (s12) [right=of s11] {\includegraphics[width=0.30\linewidth]{figures/steps_screenshot/2.jpg}};
    \node (s13) [right=of s12] {\includegraphics[width=0.30\linewidth]{figures/steps_screenshot/3.jpg}};
    \node (s21) [below=of s11] {\includegraphics[width=0.30\linewidth]{figures/steps_screenshot/4.jpg}};
    \node (s22) [right=of s21] {\includegraphics[width=0.30\linewidth]{figures/steps_screenshot/5.jpg}};
    \node (s23) [right=of s22] {\includegraphics[width=0.30\linewidth]{figures/steps_screenshot/6.jpg}};
    \node (s31) [below=of s21] {\includegraphics[width=0.30\linewidth]{figures/steps_screenshot/7.jpg}};
    \node (s32) [right=of s31] {\includegraphics[width=0.30\linewidth]{figures/steps_screenshot/8.jpg}};
    \node (s33) [right=of s32] {\includegraphics[width=0.30\linewidth]{figures/steps_screenshot/9.jpg}};
    \node (s41) [below=of s31] {\includegraphics[width=0.30\linewidth]{figures/steps_screenshot/10.jpg}};
    \node (s42) [right=of s41] {\includegraphics[width=0.30\linewidth]{figures/steps_screenshot/11.jpg}};
    \node (s43) [right=of s42] {\includegraphics[width=0.30\linewidth]{figures/steps_screenshot/12.jpg}};
    \node (l11) [draw=black,fill=white,text=black, circle, below right=0pt of s11.north west]{1};
    \node (l12) [draw=black,fill=white,text=black, circle, below right=0pt of s12.north west]{2};
    \node (l13) [draw=black,fill=white,text=black, circle, below right=0pt of s13.north west]{3};
    \node (l21) [draw=black,fill=white,text=black, circle, below right=0pt of s21.north west]{4};
    \node (l22) [draw=black,fill=white,text=black, circle, below right=0pt of s22.north west]{5};
    \node (l23) [draw=black,fill=white,text=black, circle, below right=0pt of s23.north west]{6};
    \node (l31) [draw=black,fill=white,text=black, circle, below right=0pt of s31.north west]{7};
    \node (l32) [draw=black,fill=white,text=black, circle, below right=0pt of s32.north west]{8};
    \node (l33) [draw=black,fill=white,text=black, circle, below right=0pt of s33.north west]{9};
    \node (l41) [draw=black,fill=white,text=black, circle, below right=0pt of s41.north west]{10};
    \node (l42) [draw=black,fill=white,text=black, circle, below right=0pt of s42.north west]{11};
    \node (l43) [draw=black,fill=white,text=black, circle, below right=0pt of s43.north west]{12};
    \end{tikzpicture}
    \caption{Screenshots from User Interface in Blender}
    \Description{Image shows the UI provided by the Blender and the implemented Add-On Interface. It shows 12 steps that goes through how to design channels and prepare the object for printing. Steps 1-6 shows how to designs first channel, steps 7-9 shows how to design the second channel, steps 10-11 shows how to apply partitioning and finally step 12 shows how to reorient the partitions for printing.}
    \label{fig:blenderui}
\end{figure*}



  \section{User Interface for Workflow Inside Blender\label{app:blender}}
As shown in Figure \ref{fig:blenderui} all of the design steps are implemented in Blender. Here is the summary of the steps shown in the figure:

\begin{enumerate}[noitemsep, leftmargin=*]
    \item Select object of interest in Blender viewport
    \item From the Add-On Menu, go into the \textit{Curve Partitioning Tool} and click on \verb|Add Canvas and Draw Channels|. Then select two points on the object. After selecting two points, tool will generate a plane that can be rotated with the mouse movement.
    \item Once satisfied with the plane orientation, confirm the orientation by clicking right mouse button. 
    \item Context is set to \verb|curve drawing mode| after confirming plane orientation
    \item Curves are drawn on the plane, making it easier for users to draw planar curves in 3D. At this stage all curve tools are available, whether it be deleting, drawing multiple curves, or changing vertex locations. %, all are permitted.
    \item Once satisfied with the drawn curves, they confirm the curve by pressing \verb|escape|
    \item To draw another curve, go over the same actions with clicking on \verb|Add Canvas and Draw Channels| on the Add-On menu, select points and orient the drawing plane.
    \item Draw a new curve
    \item Once satisfied with all the curves finalize the curves.
    \item Click on \verb|Apply Partitioning| on the Add-On menu to apply curve based partitioning to the object using Voronoi.
    \item Last step is to apply a procedural modifier to the object that partitions the object. There are several parameters that can be modified such as channel size, gap between the partitions, etc.
    \item Click on \verb|Apply 2D Printing Layout| on the Add-On menu to align each of the partitions and adjust the spacing between the partitions. Since all the workflow is implemented in a procedural geometry creation way, user can change parameters by going into the modifier and changing the values.
\end{enumerate}